\documentclass[a4paper,11pt]{article}
\usepackage[utf8]{inputenc}
\usepackage[T1]{fontenc}
\usepackage[english]{babel}
\usepackage{booktabs}

\begin{document}

% =====================================================================
% Data fusion part: corrected paragraphs and additional tables
% Needs: \usepackage{booktabs}, optionally \usepackage{graphicx} (for \resizebox)
% All numbers are taken from data_fusion.ipynb. "TODO" = please verify.
% =====================================================================


% ---------------------------------------------------------------------
% 1.1  Groups + preprocessing (new, short; put before "Provenance")
% ---------------------------------------------------------------------
\paragraph{Preparation.}
Before fusion, titles and author lists are cleaned per source with the
normalisation shared with identity resolution: trailing disambiguation, series and
genre notes are removed from titles (changed for 2.0\,\% of the Goodreads,
45.0\,\% of the Amazon and 23.7\,\% of the DBpedia records), and author lists are
cleaned source-specifically, e.g.\ translators and illustrators are dropped from
Goodreads lists and DBpedia disambiguations such as ``(writer)'' are removed
(55.5\,\% of the Goodreads, 9.2\,\% of the DBpedia records). The original values
are kept in \texttt{title\_raw} and \texttt{authors\_raw}. All 11,393 groups are
clean (at most one record per dataset), so filtering suspicious groups would have
no effect and all correspondences are used.


% ---------------------------------------------------------------------
% 1.1  Provenance (replaces the old paragraph)
% ---------------------------------------------------------------------
\paragraph{Provenance Metadata.}
Two provenance elements are attached to every dataset: a trust score
(\texttt{PRIOR\_TRUST}, data quality) and a dataset date
(\texttt{DATASET\_DATE}, timeliness). As we had no prior knowledge about the
reliability of the sources, all start with the same trust score of 1.0. Reliability
was instead measured per attribute on the development set as the share of correct
values among the values a source provides (Table~\ref{tab:reliability}). The
resulting source order is passed to \texttt{prefer\_higher\_trust} and to our own
\texttt{publisher\_vote} as an attribute-specific \texttt{trust\_map}, because the
sources are good at different attributes: Amazon is by far the least reliable source
for page counts (0.25), DBpedia the most reliable (0.81). For attributes without
ground truth (price, language, format, edition, genres, awards) the order follows
domain knowledge (Table~\ref{tab:strategy}).
Amazon is assigned the year 2023 and DBpedia 2026 according to the websites of
the datasets; since the collection date of Goodreads is unknown, its publication
year (2020) is used. The dates are only used by \texttt{most\_recent}, which was
tried for \texttt{publication\_date} (most book attributes are static after
publication). The original record IDs of all fused records are retained
(\texttt{\_fusion\_sources}) for traceability.

\begin{table}[t]
\centering
\caption{Share of correct values per source (tolerant comparison, development
set, only where the source has a value). Basis for the source orders.}
\label{tab:reliability}
\begin{tabular}{lccc}
\toprule
Attribute & Goodreads & Amazon & DBpedia \\
\midrule
title            & 0.71 & 0.90 & 0.89 \\
authors\_author  & 0.95 & 0.90 & 0.95 \\
publisher        & 0.85 & 0.80 & 0.84 \\
publication\_date& 1.00 & 1.00 & 0.94 \\
page\_count      & 0.70 & 0.25 & 0.81 \\
\bottomrule
\end{tabular}
\end{table}


% ---------------------------------------------------------------------
% 1.1  Density paragraph (replaces the old one; fixes Goodreads/DBpedia
%      and the unfinished singleton sentence)
% ---------------------------------------------------------------------
\paragraph{Density.}
Table~\ref{tab:density} shows the density of the input and the fused datasets; the
last row is the unweighted mean over the 10 fused attributes. Goodreads and Amazon
are dense (0.94 each), whereas DBpedia is much sparser (0.49) because it lacks
\texttt{format}, \texttt{rating}, \texttt{ratings\_count} and \texttt{price}.
The 151,748 input records are merged into 138,922 records (12,826 merged away),
i.e.\ 79,652 additional entities (+134.4\,\%) compared to the largest input dataset
(DBpedia, 59,270 records). The mean density rises from 0.492 to 0.776
(+0.284), with the largest gains in the attributes DBpedia lacks (slot filling).
It stays below that of Goodreads and Amazon because the integrated dataset also
contains the 50,724 DBpedia-only records, which keep their low density (0.486).
The 11,393 merged groups alone reach a density of 0.982, so slot filling was
effective.


% ---------------------------------------------------------------------
% 1.1  Consistency: one clarifying sentence for the last paragraph
% ---------------------------------------------------------------------
% Replace "Overall, 91.2% ..." by:
Overall, 91.2\,\% of the 11,393 merged groups contain at least one conflict across
the ten attributes under the tolerant comparison (97.0\,\% with identical values),
with on average 2.93 conflicting attributes per group. Conflict resolution is thus
relevant for almost all merged groups.


% ---------------------------------------------------------------------
% 1.2  Evaluation set: construction (replaces the old "Construction")
% ---------------------------------------------------------------------
\paragraph{Construction.}
The evaluation set contains 48 entities (238 values) for \texttt{title},
\texttt{authors\_author}, \texttt{publisher}, \texttt{publication\_date} and
\texttt{page\_count}. We exported 90 candidate groups in two batches (60 + 30),
randomly sampled from the record groups with at least two conflicting attributes;
the second batch was additionally restricted to groups whose records share an
ISBN. % TODO: verify "79 looked up" (rows in the two filled CSVs: 51 + 28)
79 candidates were looked up, and 48 of them could be verified with at least four
attributes (\texttt{page\_count} was available for 46). Eleven candidates were
dropped, e.g.\ because their group changed in the final identity-resolution run or
because no matching WorldCat edition existed. Three retained entities
changed their group after the final run: for two the ISBNs still show the same
edition, one is a wrong match and counts as identity-resolution error.
The true values were looked up manually in WorldCat via the ISBN. % TODO: keep only if true
Where an ISBN referred to several editions, the edition matching year and
publisher of the candidate was chosen.
WorldCat is an independent global library catalogue with edition-level metadata and
is more reliable than user-edited platforms such as Goodreads. The 48 entities
were split evenly into a development set (function selection) and a test set
(final evaluation), each with 24 entities and 119 values.


% ---------------------------------------------------------------------
% 1.2  Upper bound (replaces the last paragraph of 1.2; reconciles 19/17)
% ---------------------------------------------------------------------
For 19 of the 238 values (8.0\,\%) the true value does not occur in any record of
the group (12 page counts, 3 publishers, 3 titles, 1 publication year), which
bounds the accuracy of any extractive fusion rule by 92.0\,\% under tolerant
matching. One of the three titles belongs to a suspected wrong match
(Section~\ref{sec:errors}).


% ---------------------------------------------------------------------
% 1.3  Tested strategies (replaces "Tested Strategies and Custom Resolution
%      Function" AND the first sentences of "Conflict Resolution Strategies")
% ---------------------------------------------------------------------
\paragraph{Tested Strategies and Custom Function.}
For each evaluated attribute 4--6 conflict resolution functions of PyDI were
compared (Table~\ref{tab:cr}): \texttt{longest\_string}, \texttt{shortest\_string}
and \texttt{voting} for text, \texttt{union}, \texttt{intersection} and
\texttt{intersection\_k\_sources} for author lists, \texttt{earliest} and
\texttt{most\_recent} for dates, \texttt{median}, \texttt{average},
\texttt{maximum} and \texttt{minimum} for page counts, and the provenance-aware
\texttt{prefer\_higher\_trust}, which was tried with all six source orders because
reliability differs per attribute. Additionally we implemented
\texttt{publisher\_vote}, because standard voting treats spelling variants such as
``Viking'' and ``Viking Press'' as different values. Our function votes over
normalised publisher names (legal forms and generic words removed), breaks ties
by source trust and returns the most frequent spelling of the winning group.
For attributes without ground truth the functions were fixed by reasoning
(Table~\ref{tab:strategy}): median for rating and ratings\_count, trust-based
selection for price, language and format; for genres Goodreads first instead of a
union, because the genre vocabularies hardly overlap (26.1\,\% of the matched
groups share at least one genre, none has identical lists); Amazon first for
edition (more complete) and Goodreads for awards (only source).


% ---------------------------------------------------------------------
% 1.3  Table: final strategy (new)
% ---------------------------------------------------------------------
\begin{table}[t]
\centering
\caption{Final strategy. G = Goodreads, A = Amazon, D = DBpedia; the order
is the trust order of \texttt{prefer\_higher\_trust}.}
\label{tab:strategy}
\small
\begin{tabular}{llll}
\toprule
Attribute & Function & Order & Chosen by \\
\midrule
title             & shortest\_string     & --        & dev grid \\
authors\_author   & prefer\_higher\_trust& G > A > D & dev grid \\
publisher         & prefer\_higher\_trust& G > D > A & dev grid \\
publication\_date & voting               & --        & dev grid \\
page\_count       & minimum              & --        & dev grid \\
rating, ratings\_count & median          & --        & fixed \\
price, format, language & prefer\_higher\_trust & A > G > D & fixed \\
edition           & prefer\_higher\_trust& A > G > D & fixed \\
genres, awards    & prefer\_higher\_trust& G > A > D & fixed \\
\bottomrule
\end{tabular}
\end{table}


% ---------------------------------------------------------------------
% 1.3  Results text (replaces the paragraph after Table 3)
% ---------------------------------------------------------------------
On the test set the tolerant accuracy rises from 0.824 to 0.832 and the strict
accuracy from 0.681 to 0.748 (Table~\ref{tab:final}). \texttt{title},
\texttt{publication\_date} and the strict \texttt{page\_count} accuracy improve,
whereas \texttt{authors\_author} (0.958 $\to$ 0.917) and \texttt{publisher}
(0.667 $\to$ 0.625) are slightly worse than the baseline on the test set, and
tolerant \texttt{page\_count} stays at 0.652. The custom \texttt{publisher\_vote}
reached the same tolerant development accuracy as \texttt{prefer\_higher\_trust}
(0.875), but a lower strict one (0.458 vs.\ 0.583), so the simpler function was
selected.

% ---------------------------------------------------------------------
% 1.3  Table 4 extended (replaces Table 4: adds strict test accuracy)
% ---------------------------------------------------------------------
\begin{table}[t]
\centering
\caption{Accuracy of baseline (longest string, union, median) and final strategy.
Dev/test tolerant, test additionally strict. Overall is weighted by the number
of values (119 each).}
\label{tab:final}
\small
\begin{tabular}{llcccccc}
\toprule
 & & \multicolumn{2}{c}{Dev (tol.)} & \multicolumn{2}{c}{Test (tol.)} & \multicolumn{2}{c}{Test (strict)} \\
\cmidrule(lr){3-4}\cmidrule(lr){5-6}\cmidrule(lr){7-8}
Attribute & Selected function & Base. & Final & Base. & Final & Base. & Final \\
\midrule
title       & shortest string    & 0.750 & 0.875 & 0.875 & 0.958 & 0.875 & 0.958 \\
authors     & trust (G > A > D)  & 0.917 & 0.958 & 0.958 & 0.917 & 0.917 & 0.875 \\
publisher   & trust (G > D > A)  & 0.792 & 0.875 & 0.667 & 0.625 & 0.458 & 0.458 \\
pub.\ date  & voting             & 0.917 & 0.958 & 0.958 & 1.000 & 0.958 & 1.000 \\
page count  & minimum            & 0.652 & 0.783 & 0.652 & 0.652 & 0.174 & 0.435 \\
\midrule
Overall     &                    & 0.807 & 0.891 & 0.824 & 0.832 & 0.681 & 0.748 \\
\bottomrule
\end{tabular}
\end{table}

% ---------------------------------------------------------------------
% 1.3  Tolerant vs. strict (short paragraph, new)
% ---------------------------------------------------------------------
\paragraph{Strict vs.\ tolerant evaluation.}
Strict accuracy compares exactly (after trimming and lower-casing); the tolerant
one accepts title and author lists with word-Jaccard $\geq 0.8$, publishers with the
same normalised name, publication dates in the same year (WorldCat provides only
years in our evaluation set) and page counts within $\pm 2\,\%$.
The tolerance cuts both ways: it accepts ``Dune'' vs.\ ``Dune (novel)'' and the
false positive ``Penguin Books'' = ``Penguin Press'', but still rejects
``Dutton'' vs.\ ``E.\ P.\ Dutton''.


% ---------------------------------------------------------------------
% 1.3  OPTIONAL: provenance vs. no provenance (test set)
% ---------------------------------------------------------------------
\begin{table}[t]
\centering
\caption{Best function without and with provenance (chosen on dev, tolerant
accuracy on the test set).}
\label{tab:prov}
\small
\begin{tabular}{lllcc}
\toprule
Attribute & Without provenance & With provenance & Test w/o & Test with \\
\midrule
title       & shortest\_string & trust (A > D > G) & 0.958 & 0.917 \\
authors     & union            & trust (G > A > D) & 0.958 & 0.917 \\
publisher   & voting           & trust (G > D > A) & 0.708 & 0.625 \\
pub.\ date  & voting           & trust (G > A > D) & 1.000 & 1.000 \\
page count  & minimum          & trust (G > D > A) & 0.652 & 0.609 \\
\bottomrule
\end{tabular}
\end{table}
% Text: Provenance-based functions did not outperform value-only functions on the
% test set (differences of 0.04--0.08, i.e.\ one or two entities). The source
% order chosen for publisher on dev (0.875 vs. 0.833 for voting) did not
% generalise, which hints at overfitting to the small development set.
%
% NOTE: Table 4 (final) uses trust for authors/publisher, i.e. the dev-selected
% function, not the better test function. Keep this wording consistent.


% ---------------------------------------------------------------------
% 1.4  Error analysis (replaces the whole section text)
% ---------------------------------------------------------------------
\subsection{Error Analysis and Pipeline Summary}\label{sec:errors}
On the whole evaluation set (development and test, 238 values) 33 values remain
wrong or missing, a tolerant accuracy of 86.1\,\% (the functions were chosen on the
development half, so this number is slightly optimistic compared to the 83.2\,\%
on the test half). Each error was assigned one cause (Table~\ref{tab:errors}).
Most errors occur in \texttt{page\_count} (13) and \texttt{publisher} (12),
followed by \texttt{title} (4), \texttt{authors\_author} (3) and
\texttt{publication\_date} (1).

\begin{table}[t]
\centering
\caption{Remaining errors per attribute and cause (whole evaluation set).}
\label{tab:errors}
\small
\begin{tabular}{lcccc|c}
\toprule
Attribute & ID-res.\ error & No value & Wrong value chosen & True value in no source & Total \\
\midrule
title            & 1 & 0 & 1 & 2  & 4  \\
authors\_author  & 0 & 0 & 3 & 0  & 3  \\
publisher        & 1 & 0 & 8 & 3  & 12 \\
publication\_date& 0 & 1 & 0 & 0  & 1  \\
page\_count      & 0 & 0 & 1 & 12 & 13 \\
\midrule
Total            & 2 & 1 & 13 & 17 & 33 \\
\bottomrule
\end{tabular}
\end{table}

\paragraph{Not solvable by fusion (18 errors).}
In 17 cases the true value is in no source of the group, and in one case no source
has a value at all. For \texttt{page\_count} (12 cases) the sources differ by
3--7\,\% from WorldCat (e.g.\ 320 vs.\ 352, 261 vs.\ 246), caused by other printings
or counting conventions; this exceeds the 2\,\% tolerance, and no fusion function
can choose a value that does not exist. For \texttt{publisher} (3 cases) the
sources name the parent publisher or a later imprint (e.g.\ Abrams vs.\ Amulet
Books, Gallery Books vs.\ Pocket Books), which would need an external imprint
mapping. For \texttt{title} (2 cases) all sources carry the subtitle while the
reference has the main title only (``Hush: An Irish Princess' Tale'' vs.\
``Hush'').

\paragraph{Solvable in principle (13 errors).}
All eight wrong \texttt{publisher} choices are naming variants where the true
value is in a source: initials (``Dutton'' vs.\ ``E.\ P.\ Dutton''), imprint
suffixes (``Atheneum'' vs.\ ``Atheneum Books for Young Readers'', ``William
Morrow'' vs.\ ``William Morrow Paperbacks'') or merged spellings
(``HarperCollinsChildren'sBooks''); better normalisation (initials, imprint terms,
compound names like ``MTV Books/Pocket Books'') would resolve several of them.
Two ties between two sources (NAL vs.\ Penguin Books) can only be decided by an
imprint mapping. For \texttt{authors\_author}, a truncated name (``Kristen Pro'')
could be detected by a prefix check, and a subset list (James Patterson only)
by combining lists, whereas a username listed as co-author (``Fallzswimmer'') needs
the opposite rule; with only two sources there is no majority, so no single
rule fixes all three. The wrong title (``Caller'' instead of ``The Caller'') is a
side effect of \texttt{shortest\_string} that article-aware normalisation would
avoid, and the page count error (400 vs.\ 616) would be avoided by trusting DBpedia,
the most reliable source for this attribute, but \texttt{minimum} was better on the
development set overall. Such rules involve trade-offs and may introduce new errors.

\paragraph{Identity resolution (2 errors).}
For one entity the group contains records with different ISBNs, i.e.\ different
editions; an ISBN consistency check in identity resolution would have prevented
this match. The other pairs a German with an English title of the same work and
needs a clearer definition of edition-level identity.

Overall, some errors could be reduced by better normalisation, source selection
and identity resolution, whereas missing or incompatible source values (18 of 33
errors) cannot be fixed by the fusion strategy. Given the small evaluation set,
a single entity changes the accuracy by about four percentage points, so the
results should be interpreted cautiously.

\end{document}